\subsection{\texorpdfstring{ADDITION AND MULTIPLICATION IN \(\bar{R}\)}{ADDITION AND MULTIPLICATION IN \textbackslash bar\{R\}}}

Note first that the function \(-x\) can be extended by continuity to \(\overline{\mathbf{R}}\) , according to the formulae \(-(-\infty)=+\infty\) and \(-(+\infty)=-\infty\) ; the function thus extended is a homeomorphism of \(\overline{\mathbf{R}}\) onto itself.

Next, consider, the functions \(x + y\) and xy, defined on \(R \times R\) , with values in R; by considering that they take their values in the topological space \(\overline{R}\) , we shall see that they too can be extended by continuity at certain points of \(\overline{R} \times \overline{R}\) .

So far as \(x + y\) is concerned, let \(A' = ] - \infty, + \infty]\) and \(A'' = [-\infty, + \infty[\) . Then we have the following proposition:

PROPOSITION 7. The function \(x + y\) can be extended by continuity to each of the sets \(\mathbf{A}' \times \mathbf{A}'\) and \(\mathbf{A}'' \times \mathbf{A}''\) , according to the formulae

\[
\left\{ \begin{array}{l l} x + (+ \infty) = (+ \infty) + x = + \infty & (x \neq - \infty), \\ x + (- \infty) = (- \infty) + x = - \infty & (x \neq + \infty). \end{array} \right.\tag{2}
\]

Let us show, for example, that as \((x,y)\) tends to the point \((a,+\infty)\) \((a\neq-\infty)\) while remaining in \(R\times R\) , \(x+y\) tends to \(+\infty\) . There exists a finite number b\textless a, and the interval {]}b, \(+\infty\) {[} is a neighbourhood of a in R; given any finite c, the relations x\textgreater b and y\textgreater c-b imply \(x+y>c\) , and this shows that \(x+y\) is as near as we please to \(+\infty\) whenever \((x,y)\) is near enough to \((a,+\infty)\) . The argument is similar in the other cases.

On the contrary, \(x + y\) has no limit at the points \((-\infty, +\infty)\) and \((+\infty, -\infty)\) of \(\overline{\mathbf{R}} \times \overline{\mathbf{R}}\) . For if \(x + y\) had limit \(k\) (finite or infinite) as \((x, y)\) tends to \((+\infty, -\infty)\) while remaining in \(\mathbf{R} \times \mathbf{R}\) , it would follow that, for each finite \(a\) , the function \((x + a) - x\) would tend to \(k\) as \(x\) tends to \(+\infty\) while remaining in \(\mathbf{R}\) ; and this is absurd, since \((x + a) - x = a\) and \(a\) is arbitrary.

The function \(x + y\) maps \(A' \times A'\) (resp. \(A'' \times A''\) ) into \(A'\) (resp. \(A''\) ). It is therefore a law of composition on \(A'\) (resp. \(A''\) ) which extends the law of addition on \(R\) . By the principle of extension of identities (Chapter I, § 8, no. 1, Corollary 1 to Proposition 2) this law is commutative and associative; o is the identity element for this law; and the only non-regular element (Algebra, Chapter I, § 2, no. 2) of \(A'\) is \(+\infty\) , by formulae (2).

If \(x, y, z, t\) are points of \(\overline{\mathbf{R}}\) such that \(x \leqslant y\) and \(z \leqslant t\) , then \(x + z \leqslant y + t\) whenever both sides of this inequality are defined.

Note that in \(\overline{\mathbf{R}}\) the relation \(x < y\) implies \(x + z < y + z\) only if \(z\) is finite, by formulae (2); it is easily verified that the relations \(x < y\) and \(z < t\) imply \(x + z < y + t\) whenever both sides of this inequality are defined.

In \(\overline{\mathbf{R}}\) we put \(x^{+} = \sup (x,0)\) , \(x^{-} = \sup (-x,0)\) , \(|x| = \sup (x, - x)\) ; thus \((+\infty)^{+} = (-\infty)^{-} = +\infty\) and \((+\infty)^{-} = (-\infty)^{+} = 0\) , and \(|+\infty| = |- \infty| = +\infty\) . The sums \(x^{+} - x^{-}\) and \(x^{+} + x^{-}\) are defined for all \(x \in \overline{\mathbf{R}}\) and are therefore equal to \(x\) and \(|x|\) respectively by the principle of extension of identities. Also, whenever the sum \(x + y\) is defined, we have \(|x + y| \leqslant |x| + |y|\) .

Note on the contrary that formulae (6) and (7) of § 1, no. 1 may no longer make sense for certain values of \(x\) and \(y\) in \(\overline{\mathbf{R}}\) ; for example, if \(x = -\infty\) and \(y = 0\) we have \(\sup (x, y) = 0\) , but the sum \(x + (y - x)^+\) is not defined, because \((y - x)^+ = +\infty\) .

Let \(\overline{\mathbf{R}}^{*}\) denote the complement of \(o\) in \(\overline{\mathbf{R}}\) . Then the analogue of Proposition 7 for multiplication runs as follows:

PROPOSITION 8. The function \(xy\) can be extended by continuity to the set \(\overline{\mathbf{R}}^* \times \overline{\mathbf{R}}^*\) according to the formulae

\[
\left\{ \begin{array}{l} x. (+ \infty) = (+ \infty). x = \left\{ \begin{array}{l l} + \infty & \text {if} \quad x > 0 \\ - \infty & \text {if} \quad x <   0 \end{array} \right. \\ x. (- \infty) = (- \infty). x = \left\{ \begin{array}{l l} - \infty & \text {if} \quad x > 0 \\ + \infty & \text {if} \quad x <   0. \end{array} \right. \end{array} \right.\tag{3}
\]

We leave the proof to the reader; it is analogous to the proof of Proposition 7.

Likewise, we see that \(xy\) has no limit at the points \((0, +\infty), (+\infty, 0), (0, -\infty), (-\infty, 0)\) of \(\overline{\mathbf{R}} \times \overline{\mathbf{R}}\) .

The function \(xy\) is a law of composition on \(\overline{\mathbf{R}}^*\) which extends the law of multiplication on \(\mathbf{R}\) ; this law is associative and commutative (principle of extension of identities); it has 1 as identity element; and the non-regular elements in \(\overline{\mathbf{R}}^*\) are \(+\infty\) and \(-\infty\) .

If \(x \leqslant y\) and \(z > 0\) we have \(xz \leqslant yz\) whenever both sides of this inequality are defined. If the product \(xy\) is defined, then so is \(|x| \cdot |y|\) and we have \(|xy| = |x| \cdot |y|\) .

Finally, the distributivity formula

\[
x (y + z) = x y + x z\tag{4}
\]

is still valid, by virtue of the principle of extension of identities, whenever all the operations which figure on either side are defined.

\pandocbounded{\includegraphics[keepaspectratio]{images/part002_cca5eabe.jpg}}

Note that it can happen that the left-hand side of (4) is defined while the right-hand side is not: for example, consider the case where \(x = +\infty\) , y = 2 and z = -1. The distributivity formula should therefore be used with caution in R.

